\section*{Reproducibility and Artifact Access}
\addcontentsline{toc}{section}{Reproducibility and Artifact Access}
The report follows the handbook's three-part structure and distinguishes measured
results from proposed follow-ups.\citep{handbook}
The source repository is
\href{https://github.com/TuanKhai1210/Deep-Learning-Benchmark-Suite}{Deep-Learning-Benchmark-Suite};
the assignment entry point is
\href{https://tuankhai1210.github.io/Deep-Learning-Benchmark-Suite/a1.html}{the A1 project page}.
The reported run revision is:
\begin{center}\small\texttt{a8539bee778d99e2902a1dd36aa157c6ebb3a42c}\end{center}
The run artifacts and snapshots, rather than an older README status paragraph,
define the effective configurations. Appendix~\ref{app:runs} maps all 15 results
to their run IDs; the source package includes exact effective JSON configs,
per-run metrics, raw timing samples and input/checkpoint hashes under
\texttt{evidence/}.

Original checkpoints are maintained in the \href{https://drive.google.com/drive/u/1/folders/13QiG0sOMNTxxGe8er4BesVr4Aky4IL73}{group's Drive artifact bundle}, with one run directory per ID. Readers without that bundle can reconstruct the models by training the included effective configs at the pinned source revision. A copy of the checkpoints are available on the source repository.

For a Linux/Colab CUDA environment, the following installation identifies the
main recorded tensor-library builds. Python patch versions and operational
differences remain as described in Section~\ref{sec:environment}.
\begin{lstlisting}[basicstyle=\ttfamily\footnotesize,breaklines=true,frame=single]
git clone https://github.com/TuanKhai1210/Deep-Learning-Benchmark-Suite.git
cd Deep-Learning-Benchmark-Suite
git checkout a8539bee778d99e2902a1dd36aa157c6ebb3a42c
python -m pip install torch==2.11.0 torchvision==0.26.0 --index-url https://download.pytorch.org/whl/cu130
python -m pip install numpy==2.1.3
python -m pip install -e ".[ml]"
\end{lstlisting}
Official previous-version installation instructions document these builds.\citep{pytorchversions}
The effective configs and the environment records govern reconstruction;
installing unpinned optional extras alone is not a complete environment lock.

The included \path{reproduction/reproduce.py} reads a selected saved config,
requires the pinned repository revision and changes only explicit operational
paths/device settings. Its default action prints the resolved configuration;
training requires the explicit action below. In this example the report
source archive has been extracted to \texttt{/content/A1\_Final}:
\begin{lstlisting}[basicstyle=\ttfamily\footnotesize,breaklines=true,frame=single]
python /content/A1_Final/reproduction/reproduce.py \
  --repo-root . --model cnn --seed 67 --action train \
  --data-root data --output-root runs/a1_reproduction --device cuda
\end{lstlisting}
Repeat with the five model names \texttt{linear, mlp, cnn, rnn, transformer}
and seeds 69420, 67, 69. A new run may differ numerically because exact cross-machine
or cross-version bitwise reproduction has not been established.
Evaluation of a trusted saved checkpoint requires an explicit action and path:
\begin{lstlisting}[basicstyle=\ttfamily\footnotesize,breaklines=true,frame=single]
python /content/A1_Final/reproduction/reproduce.py \
  --repo-root . --model cnn --seed 67 --action test \
  --checkpoint /content/drive/MyDrive/dlbench_a1_main/main/a1_cnn_main_seed67_20261023-030133_702c4e/best.pt \
  --data-root data --device cpu
\end{lstlisting}
The source package also retains the one-pass evaluation notebook and timing-only
notebook in \texttt{reproduction/}. Their outputs preserve split/device labels
and completion identities. Regenerating report tables and figures from existing
artifacts uses \path{tools/build_assets.py}; it performs no model training or
checkpoint inference.

\section*{Member Contributions and AI Usage Disclosure}
\addcontentsline{toc}{section}{Member Contributions and AI Usage Disclosure}
The task division follows the project's recorded ownership: Nguyen Hao Thien
is responsible for data/EDA/preprocessing and Linear/CNN; Nguyen Anh Khoa for
the shared training/evaluation engine, MLP/GRU and metrics/checkpointing;
Ta Tuan Khai for Transformer, configuration/artifacts, integration and report
reproducibility. Shared interfaces and results require cross-member review.
The cover sheet retains the group's supplied completion declarations; automated
artifact checks do not independently measure individual contribution percentages.

AI assistance was used and is disclosed below. Historical details come from
the repository's \texttt{AI\_USAGE.md}, copied into the source package as
\path{evidence/AI_USAGE_source.md}. Model/version labels in that log are
reported as recorded rather than independently certified service identifiers.
\begin{table}[H]\centering\small
\caption{Recorded AI tools, users and development scope.}
\label{tab:ai}
\begin{tabularx}{\linewidth}{>{\raggedright\arraybackslash}p{30mm}>{\raggedright\arraybackslash}p{29mm}>{\raggedright\arraybackslash}X}
\toprule Tool/model as recorded & Member & Stage and affected work\\\midrule
OpenAI Codex & Ta Tuan Khai & Aug.--Oct.\ 2026: protocol, artifacts, Transformer support,
analysis, source correction and report generation\\
Gemini & Nguyen Hao Thien; Nguyen Anh Khoa & Sep.\ 2026: data/EDA, metrics/checkpoints,
trainer tests and analysis notebook\\
GitHub Copilot & Nguyen Anh Khoa & Sep.\ 2026: MLP/trainer tests and optimizer support\\
OpenAI ChatGPT & Nguyen Anh Khoa & Sep.\ 2026: trainer test-generation support\\
\bottomrule
\end{tabularx}
\end{table}

The historical log names Codex 5.6 Sol and GPT-6 in selected entries, Gemini
3.8 Flash and 3.1 Pro (High), and ChatGPT 5.6 Luna. Other serving model IDs
were not recorded. Representative prompt summaries include explaining leakage-free EDA and stratified
splits (Thien), identifying states needed for training resumption and generating
metric/trainer tests (Khoa), and making artifacts traceable to configs/checkpoints
and completing the LaTeX report from verified runs (Khai).
Historical log entries record student checks against the handbook, source interfaces,
PyTorch documentation and tests. They do not establish that AI output is correct
without those checks.

For the 23--26 October report/analysis work, Khai used Codex; the exact serving
model identifier for every interaction was not retained. AI proposed and edited
report prose, generated aggregation/plotting helpers, corrected source/report
inconsistencies, and helped assemble the final document. The current automated
verification recomputes aggregates from saved predictions/metrics, checks source
and checkpoint identities, verifies raw dataset resources, and compiles/inspects
the PDF. Official documentation and papers support the theoretical descriptions.
The team remains responsible for the interpretation and final verification of
its submission; this disclosure does not assert an unrecorded human sign-off.
Additional session detail is retained in \path{evidence/AI_USAGE_report.md}.
\FloatBarrier

\clearpage
\begin{thebibliography}{9}
\bibitem{fashion}
H. Xiao, K. Rasul and R. Vollgraf.
\emph{Fashion-MNIST: a Novel Image Dataset for Benchmarking Machine Learning Algorithms}.
2017. \url{https://arxiv.org/abs/1708.07747}.
\bibitem{fashionlicense}
Zalando Research. \emph{Fashion-MNIST repository and MIT License}.
2017. \url{https://github.com/zalandoresearch/fashion-mnist/blob/master/LICENSE}.
Accessed 26 October 2026.
\bibitem{attention}
A. Vaswani et al. \emph{Attention Is All You Need}.
Advances in Neural Information Processing Systems, 2017.
\url{https://arxiv.org/abs/1706.03762}.
\bibitem{pytorchgru}
PyTorch contributors. \emph{torch.nn.GRU, PyTorch 2.11 documentation}.
\url{https://docs.pytorch.org/docs/2.11/generated/torch.nn.GRU.html}.
Accessed 26 October 2026.
\bibitem{pytorchce}
PyTorch contributors. \emph{CrossEntropyLoss, PyTorch 2.11 documentation}.
\url{https://docs.pytorch.org/docs/2.11/generated/torch.nn.CrossEntropyLoss.html}.
Accessed 26 October 2026.
\bibitem{pytorchversions}
PyTorch contributors. \emph{Installing previous versions of PyTorch}.
\url{https://pytorch.org/get-started/previous-versions/}. Accessed 26 October 2026.
\bibitem{project}
Tiki-taka. \emph{Deep Learning Benchmark Suite}, source revision a8539be.
\url{https://github.com/TuanKhai1210/Deep-Learning-Benchmark-Suite}.
\bibitem{handbook}
CO3133, Semester 261. \emph{Course Project Handbook}, revision 14 September 2026.
Course-provided document; Sections 3--5 and 8--14.
\end{thebibliography}
